\documentclass[../../../include/open-logic-section]{subfiles}

\begin{document}

\olfileid{sth}{choice}{justifications}
\olsection{પસંદગી અંગે આંતરિક વિચારણાઓ}

તો વ્યાપક પ્રશ્ન એ છે કે સુક્રમિતતા,
અથવા પસંદગી, અથવા ખરેખર કદની દૃષ્ટિએ
બધા ગણોની તુલનીયતા---આમાંથી કયું
કહીએ તેનાથી ફરક પડતો નથી---વાજબી
ઠરાવી શકાય છે કે નહીં.

આ એક \emph{આંતરિક} સમર્થનનો પ્રયાસ છે.
\olref[z][story]{sec}માં આપણે પદાનુક્રમ
વિશે કેટલાક સિદ્ધાંતો રજૂ કર્યા હતા.
એમાંથી એક ફરી કહેવા જેવો છે:
\begin{enumerate}
	\item[] \stagesacc. કોઈ પણ તબક્કા $S$
	અને તબક્કા $S$ની \emph{પહેલાં} રચાયેલા
	ગણો માટે: બરાબર એ જ ગણો ઘટકો
	તરીકે ધરાવતો ગણ તબક્કા $S$ પર
	રચાય છે. તબક્કા~$S$ પર બીજું કંઈ
	રચાતું નથી.
\end{enumerate}
ખરેખર ઘણા લેખકોએ સૂચવ્યું છે કે
આ સિદ્ધાંત (અથવા તેના જેવા કોઈ
સિદ્ધાંત) વડે પસંદગીની સ્વયંસિદ્ધિને
વાજબી ઠરાવી શકાય. એ અભિગમની
ટૂંકી સમજૂતી આપીએ.

શરૂઆતમાં એક નાનું સરળ પરિણામ લઈએ,
જે પસંદગીનું \emph{હજી એક} સમકક્ષ
સ્વરૂપ આપે છે:

\begin{thm}[$\ZF$માં]\ollabel{choiceset}
પસંદગી નીચેના સિદ્ધાંતને સમકક્ષ છે.
જો $A$ના !!{element}s પરસ્પર
વિચ્છિન્ન અને અરિક્ત હોય, તો એવો
$C$ છે કે દરેક $x \in A$ માટે
$C \cap x$ એક-ઘટકી ગણ છે.
(આવા $C$ને $A$ માટેનો
{પસંદગી ગણ} કહીએ.)
\end{thm}

આ પરિણામનો પુરાવો સીધો છે;
વાચક માટે કસરત તરીકે છોડીએ.

\begin{prob}
\olref[sth][choice][justifications]{choiceset}
સાબિત કરો. મુશ્કેલી પડે તો
\cite[pp.~242--3]{Potter2004}માં
પુરાવો મળશે.
\end{prob}

મુખ્ય વાત એ છે કે $A$ માટેનો પસંદગી
ગણ, $A$ માટેના પસંદગી વિધેયનું
ચિત્ર જ છે. તેથી પસંદગીને વાજબી
ઠરાવવા, પસંદગી ગણોના અસ્તિત્વની
તેની સમકક્ષ રજૂઆતને જ વાજબી
ઠરાવવાનો પ્રયાસ કરી શકીએ.
હવે એ જ કરવાનો પ્રયત્ન કરીએ.

$A$ના !!{element}s પરસ્પર વિચ્છિન્ન
અને અરિક્ત ધારો. \stageshier{} મુજબ
(\olref[z][story]{sec} જુઓ), $A$
કોઈ તબક્કા~$S$ પર રચાય છે.
નોંધો કે $\bigcup A$ના બધા
!!{element}s તબક્કા $S$ પહેલાં
ઉપલબ્ધ છે. હવે \stagesacc{} મુજબ,
$S$ પહેલાં રચાયેલા \emph{ગમે તે}
ગણોને બરાબર ઘટકો તરીકે ધરાવતો
ગણ રચાય છે. બીજા શબ્દોમાં,
અગાઉ ઉપલબ્ધ ગણોના દરેક
\emph{શક્ય} સંગ્રહનું $S$ પર
અસ્તિત્વ હશે. પરંતુ $A$ માટેનો
પસંદગી ગણ બની શકે એવી વસ્તુઓ
પસંદ કરવી ચોક્કસપણે \emph{શક્ય}
છે; તે $\bigcup A$નો એક ખાસ
ઉપગણ જ છે. તેથી જરૂરી પસંદગી
ગણ અસ્તિત્વમાં છે.

સારું, આ તો આંતરિક આધારે પસંદગીને
વાજબી ઠરાવવાનો \emph{ખૂબ} ટૂંકો
પ્રયાસ થયો. આ વિચારને આગળ
લઈ જવા Potterનું
(\citeyear[\S14.8]{Potter2004})
સરસ વિસ્તરણ વાંચો.

\end{document}
